%% Per-limb WiFi CSI separation without limb identity
%% IEEE conference format (IEEEtran). Compile:
%%   pdflatex phase2_ieee && bibtex phase2_ieee && pdflatex phase2_ieee x2
%% A self-contained thebibliography is also included at the end; to use the
%% .bib instead, comment that block and uncomment the \bibliography lines.
\documentclass[conference]{IEEEtran}
\IEEEoverridecommandlockouts

\usepackage{cite}
\usepackage{amsmath,amssymb,amsfonts}
\usepackage{algorithmic}
\usepackage{graphicx}
\usepackage{textcomp}
\usepackage{xcolor}
\usepackage{booktabs}
\usepackage{url}

\def\BibTeX{{\rm B\kern-.05em{\sc i\kern-.025em b}\kern-.08em
    T\kern-.1667em\lower.7ex\hbox{E}\kern-.125emX}}

%% Marks the slots deliberately left empty until the results pass.
\newcommand{\resultslot}[1]{\textcolor{red}{[\,#1\,]}}

\begin{document}

\title{Per-Limb Separation of WiFi Channel State Information\\
Without Limb Identity}

\author{\IEEEauthorblockN{Author Name}
\IEEEauthorblockA{\textit{Affiliation} \\
\textit{Institution}\\
City, Country \\
email@domain}
}

\maketitle

\begin{abstract}
Commodity WiFi sensing has progressed from coarse activity recognition to
full-body pose estimation, but decomposing the received channel into
\emph{per-limb} components remains open. We separate two questions that prior
work conflates: whether a room-stable spatial signature exists that
\emph{identifies} which limb produced a component, and whether the signal can
be \emph{partitioned} into limb-coherent components without ever naming them.
For the first question we build three instruments of increasing power on a
five-room, three-receiver dataset: a non-parametric subcarrier-by-antenna
signature, a coherent MUSIC estimator over a cross-antenna representation that
preserves angle-of-arrival, and a Doppler-gated joint estimator following the
multi-dimensional tracking literature. For the second, we adopt the
deep-clustering premise from speech separation---grouping time--frequency bins
requires no source identity, with the permutation resolved afterwards---and
build a full pipeline: a tokeniser that reduces each recording to a set of
time--frequency tokens carrying Doppler, angle, and delay coordinates, obtained
with Slepian multitaper analysis and an exact fast MUSIC evaluation; a
zero-learning clustering test with a matched-permutation null; and a set
transformer trained with two identity-free objectives, an
envelope permutation-invariant loss against keypoint-derived limb references
and a mixture-invariant origin loss requiring no labels at all. Models are
trained on three rooms and gated on two rooms never seen, always against a
paired zero-learning control. \resultslot{results summary}
\end{abstract}

\begin{IEEEkeywords}
WiFi sensing, channel state information, source separation, deep clustering,
micro-Doppler, MUSIC, multitaper, human pose estimation
\end{IEEEkeywords}

\section{Introduction}

WiFi channel state information (CSI) reacts to human motion, and a decade of
work has turned that reaction into increasingly fine-grained perception: from
activity classification~\cite{wang2015carm} and through-wall
localisation~\cite{adib2013seethrough}, to gesture recognition that transfers
across environments~\cite{zheng2019widar3}, to dense body
perception~\cite{wang2019personinwifi,geng2022densepose} and multi-person 3D
pose~\cite{yan2024personinwifi3d,perceptalign2026}. Each step has narrowed the
spatial scale at which the signal is interpreted.

The natural limit of that progression is a \emph{per-limb} decomposition: given
the CSI recorded while a person moves, produce one component stream per moving
body part. Such a decomposition would be the WiFi analogue of source separation
in audio, and it would provide a representation on which downstream tasks could
be built without re-learning the physics each time.

The obstacle is usually stated as a resolution problem. A wrist is a weak
scatterer a few tens of centimetres from the torso; a commodity receiver with
three antennas provides two coherent baselines, whose angular resolving power
is far coarser than the angle a limb subtends at typical indoor ranges. On that
reading, per-limb decomposition awaits better hardware.

We argue the problem statement hides two different questions:

\begin{enumerate}
\item \textbf{Identification.} Does there exist a signature---in the
subcarrier, antenna, angle, or delay domain---that is stable across rooms and
says \emph{which} limb produced a given component?
\item \textbf{Separation.} Can the signal be partitioned into limb-coherent
components \emph{without} such a signature, with the correspondence between
components and limbs resolved only afterwards?
\end{enumerate}

Speech separation settled the analogous distinction. Deep
clustering~\cite{hershey2016deep} trains embeddings of time--frequency bins
under an objective that says only ``bins of the same source should be close'',
never which speaker a source is; at test time the bins are clustered per
mixture and the outputs matched to references afterwards. The method separates
speakers it has never heard, and its successors keep the
principle~\cite{zeghidour2021wavesplit}. Two same-pitch talkers have no stable
identifying signature either, yet they separate---because within one mixture
their bins differ in local structure. Identification is sufficient for
separation, not necessary for it.

This paper takes both questions in turn on a five-room, three-receiver dataset
with synchronised 3D keypoints~\cite{perceptalign2026}. The contributions are
methodological:

\begin{itemize}
\item A precise account (Sec.~\ref{sec:model}) of what the commonly used
\emph{island} CSI representation---products of adjacent subcarriers within an
antenna---discards, and a cross-antenna alternative that cancels the same
per-packet impairments while preserving angle-of-arrival and the full delay
aperture.
\item A normalisation rule for the spectral domain (Sec.~\ref{sec:norm}):
because multipath superposes additively in the channel frequency response, the
static room must be \emph{subtracted}, and only its \emph{magnitude} divided
out, so that the phase ramps encoding angle and delay survive normalisation.
\item Three identification probes of increasing power (Sec.~\ref{sec:id}),
ending with a Doppler-gated joint (angle, delay, Doppler) estimator in the
spirit of mD-Track~\cite{xie2019mdtrack} and Widar2.0~\cite{qian2018widar2},
each with matched nulls and intensity controls.
\item An identity-free separation pipeline (Sec.~\ref{sec:sep}): a token
representation of each recording obtained by Slepian multitaper
analysis~\cite{thomson1982multitaper,slepian1978dpss} and per-bin
MUSIC~\cite{schmidt1986music} with an exact fast evaluation; a zero-learning
clustering feasibility test with a matched-permutation null; and a set
transformer trained by an envelope permutation-invariant loss and a
mixture-invariant~\cite{wisdom2020mixit} origin loss.
\item A transfer protocol (Sec.~\ref{sec:transfer}) that asks whether the
separator's frozen output is a universal representation, by training a small
pose head on it and testing in an unseen room against both a raw-feature arm
and a mean-pose floor.
\end{itemize}

Every reported test number is measured in rooms excluded from all fitting, and
every learned number is reported beside a zero-learning control evaluated on
the same recordings.

\section{Related Work}

\subsection{WiFi sensing and body perception}

CARM~\cite{wang2015carm} related CSI variation to body motion through a
speed--frequency model, establishing micro-Doppler as the natural domain for
activity sensing. Widar~\cite{qian2017widar} and Widar2.0~\cite{qian2018widar2}
turned Doppler into passive tracking, and Widar3.0~\cite{zheng2019widar3}
introduced body-coordinate velocity profiles as a domain-independent feature,
recognising that raw CSI statistics are tied to deployment geometry. Person-in-
WiFi~\cite{wang2019personinwifi}, DensePose-from-WiFi~\cite{geng2022densepose},
and multi-person 3D pose~\cite{yan2024personinwifi3d} regress body
representations directly. The dataset used here~\cite{perceptalign2026}
addresses cross-layout generalisation by conditioning on calibrated transceiver
geometry. None of this line produces a per-component decomposition of the
signal itself; the body is recovered as a regression target rather than as
separated sources.

\subsection{Multipath parameter estimation}

SpotFi~\cite{kotaru2015spotfi} applied MUSIC~\cite{schmidt1986music} with
spatial smoothing~\cite{shan1985smoothing} to joint angle--delay estimation on
commodity cards. Dynamic-MUSIC~\cite{li2016dynamicmusic} and
IndoTrack~\cite{li2017indotrack} exploit the Doppler dimension to isolate paths
reflected by a moving body from the static background, and
mD-Track~\cite{xie2019mdtrack} makes the argument general: paths that are
unresolvable in any single dimension become separable in the joint
(angle, delay, Doppler) space, because ambiguity in one dimension is broken by
another. Our Sec.~\ref{sec:dopplergated} applies exactly this correction to our
own Sec.~\ref{sec:coherent} estimator, and Sec.~\ref{sec:token} carries the
joint-space estimate into a per-bin tokeniser.

Cancelling the per-packet carrier and sampling frequency offsets of commodity
hardware~\cite{halperin2011tool} is a prerequisite for any coherent processing.
Two standard devices exist: differencing along the subcarrier axis within one
antenna, and taking ratios or conjugate products \emph{across} antennas of one
receiver~\cite{li2017indotrack,zeng2019farsense}. Both cancel the offsets
exactly, since they are common to all subcarriers and all antennas of a
receiver; they differ in what else they cancel, which Sec.~\ref{sec:islands}
makes explicit.

\subsection{Separation without identity}

DUET~\cite{yilmaz2004duet} separates audio mixtures by clustering
time--frequency bins on their inter-channel delay and attenuation, assuming
approximate disjointness of sources in the time--frequency plane; the cluster
identities are arbitrary. Deep clustering~\cite{hershey2016deep} replaces the
hand-designed bin coordinates with learned embeddings under a
class-independent objective, and Wavesplit~\cite{zeghidour2021wavesplit}
resolves the permutation problem by clustering per-frame speaker
representations. Permutation-invariant training~\cite{yu2017pit} makes the same
idea a loss, and mask-based models operating in the time--frequency
domain~\cite{wang2023tfgridnet} or on learned
bases~\cite{luo2019convtasnet} provide the architectures.

When isolated sources are unavailable, mixture-invariant
training~\cite{wisdom2020mixit} learns separation from mixtures alone by
requiring that estimated sources can be re-grouped to reconstruct each
constituent mixture; RemixIT~\cite{tzinis2022remixit} extends the family to
continual self-training. This matters here because isolated single-limb
recordings do not exist and cannot be collected without instrumenting the body:
supervised separation targets are unavailable \emph{in principle}, not merely
in practice.

\section{Signal Model and Representation}
\label{sec:model}

\subsection{Measurement model}

For a receiver with $A$ antennas and $S$ subcarriers, the measured CSI at
packet time $t$ is
\begin{equation}
\hat{H}_{a,s}(t) = g(t)\, e^{\,j(\alpha(t) + \beta(t)s + \theta_a)}
\, H_{a,s}(t) + n_{a,s}(t),
\label{eq:measmodel}
\end{equation}
where $g(t)$ is the automatic gain, $\alpha(t)$ the carrier frequency offset,
$\beta(t)$ the sampling frequency offset slope, $\theta_a$ the per-antenna
phase (carrying angle of arrival), and $H$ the propagation channel. Crucially,
$g$, $\alpha$, and $\beta$ are properties of the receiver chain and are
\emph{common to all antennas} of one card~\cite{halperin2011tool}, since a
single oscillator and sampling clock serve them.

The propagation channel superposes paths,
\begin{equation}
H_{a,s}(t) = \underbrace{\textstyle\sum_{k\in\mathcal{S}} \gamma_k
e^{-j2\pi f_s \tau_k} e^{-j\psi_{a,k}}}_{\text{static: room}}
+ \underbrace{\textstyle\sum_{k\in\mathcal{D}(t)} \gamma_k(t)
e^{-j2\pi f_s \tau_k(t)} e^{-j\psi_{a,k}(t)}}_{\text{dynamic: person}},
\label{eq:paths}
\end{equation}
so the room and the person add rather than multiply---a fact that dictates the
normalisation rule of Sec.~\ref{sec:norm}.

\subsection{The island representation and what it discards}
\label{sec:islands}

A widely used sanitisation forms products of \emph{adjacent subcarriers within
one antenna},
\begin{equation}
z_{a,s}(t) = \hat{H}_{a,s}(t)\,\hat{H}^{*}_{a,s+1}(t),
\end{equation}
which we call the \emph{island} representation. Substituting
\eqref{eq:measmodel}, the common terms cancel: $\alpha(t)$ exactly, and
$\beta(t)s$ down to a constant removed by a subsequent detrend. The resulting
stream is temporally coherent and is what earlier stages of this project used.

The same difference, however, cancels every quantity that is constant across
the subcarrier index of one antenna---including $\theta_a$. No pair of entries
in the island stream relates antenna $a$ to antenna $b$, so angle of arrival is
not merely attenuated but \emph{absent}; and the SFO detrend additionally
flattens the delay ramp. Cleaning and normalisation are not responsible for
this loss: it is a property of the representation. This observation motivates
the alternative below and explains why identification probes on island streams
(Sec.~\ref{sec:sigprobe}) cannot be read as evidence about the spatial domain.

\subsection{Coherent cross-antenna representation}

We form instead the conjugate products across antennas of one receiver,
\begin{equation}
y_{a,s}(t) = \hat{H}_{a,s}(t)\,\hat{H}^{*}_{1,s}(t),\qquad a = 2,\dots,A,
\label{eq:coherent}
\end{equation}
following the CSI-ratio family~\cite{li2017indotrack,zeng2019farsense}. Since
$g$, $\alpha$, and $\beta$ are common across antennas, they cancel exactly, as
in the island form; but $\theta_a - \theta_1$ survives, so angle of arrival
remains observable, and no subcarriers are consumed by the operation, so the
full delay aperture is retained. With $A = 3$ and $S = 57$ this yields a
$(T, 2, 57)$ complex stream per receiver: two coherent baselines by
fifty-seven subcarriers.

\subsection{Cleaning}
\label{sec:clean}

The following is applied whenever \eqref{eq:coherent} is built.

\emph{Rate recovery.} Packet rate is derived from timestamps, never assumed:
the median timestamp increment is tested against unit scales
$\{1, 10^{-3}, 10^{-6}, 10^{-9}\}$\,s and the scale placing the rate in
$[100, 5000]$\,Hz adopted. Non-monotonic samples are dropped and recordings
shorter than \SI{2}{\second} rejected.

\emph{Gain removal.} Each packet is divided by
$g(t) = \sqrt{\mathrm{mean}_{a,s}|\hat H_{a,s}(t)|^2}$, a real positive scalar,
so phases are untouched. This precedes \eqref{eq:coherent}, which would
otherwise square any residual gain variation.

\emph{Resampling.} The complex products are binned onto a uniform
\SI{400}{\hertz} grid; empty bins are nearest-neighbour filled and a recording
is rejected if more than \SI{35}{\percent} of bins are empty.

\subsection{Normalisation in the spectral domain}
\label{sec:norm}

By \eqref{eq:paths} the static room enters additively. The normalisation must
therefore \emph{subtract} it. Division by the complex static---the
cepstral-mean analogue borrowed from speech~\cite{chen2018cepstral} and used in
the earlier phase of this work---does remove the room and equalise gain, but it
rotates every element by $-\arg S(e)$, destroying precisely the phase ramps
across subcarriers and antennas that encode delay and angle. We therefore use
the phase-safe form: with $e = (a, s)$ indexing the coherent aperture,
\begin{align}
S(e) &= \tfrac{1}{T}\textstyle\sum_t y_e(t), \\
g(e) &= \max\!\big(|S(e)|,\; 0.05\cdot\mathrm{median}_e |S(e)|\big), \\
\tilde D_e(t) &= \big(y_e(t) - S(e)\big) \,/\, g(e).
\label{eq:phasesafe}
\end{align}
Subtraction removes the room exactly and leaves the dynamic term's complex
structure intact; the division is by a real positive scalar per element, so it
equalises strong and weak aperture elements---conditioning the covariance
estimates of Sec.~\ref{sec:coherent}---without rotating anything. The floor at
\SI{5}{\percent} of the median magnitude prevents dead elements from dominating
the quotient. The static is estimated from the recording itself, which keeps
the procedure single-recording and deployment-honest at the cost of a disclosed
limitation: a person who never moves is absorbed into $S$.

Subtraction matters even though the analysis band begins at \SI{2}{\hertz}: the
static term exceeds the dynamic term by orders of magnitude, and its leakage
through window sidelobes would otherwise flood the low-Doppler bins.

\subsection{Time--frequency analysis}
\label{sec:tf}

All spectra use \num{256}-sample windows with \num{128}-sample hop on the
\SI{400}{\hertz} grid (\SI{0.64}{\second} windows, \SI{0.32}{\second} hop),
retaining bins in $\pm 2$--\SI{150}{\hertz}. Sections~\ref{sec:sigprobe}--%
\ref{sec:dopplergated} use a Hann window. From Sec.~\ref{sec:token} onward we
use Thomson multitaper analysis~\cite{thomson1982multitaper} with $K = 4$
discrete prolate spheroidal (Slepian) sequences~\cite{slepian1978dpss} at
$NW = 2.5$. Slepian tapers maximise in-band energy concentration and are the
standard remedy where the sidelobes of a dominant component would obscure a
weaker superimposed one---the torso-versus-limb situation exactly. As a second
benefit the $K$ tapers provide $K$ nearly independent spectral estimates per
window, which serve directly as additional snapshots for the covariance
estimates of Sec.~\ref{sec:token}.

\subsection{Limb reference signals}
\label{sec:gt}

The dataset provides no inertial sensors; per-limb references derive from its
multi-view 3D keypoints (BODY25~\cite{cao2021openpose}, reconstructed with
EasyMocap~\cite{easymocap}). For five joints---left wrist, right wrist, left
hip, right hip, head---keypoints with confidence above \num{0.3} are converted
to 3D speeds, resampled to the CSI grid, smoothed with a \SI{0.5}{\second}
moving average, and normalised to unit variance.

Alignment follows the dataset's own convention, applied identically to all five
rooms: keypoint frames and CSI are co-started and mapped proportionally, then
interpolated onto the CSI time base. The released preprocessing likewise maps
each keypoint frame to a proportional CSI segment, with no per-frame timestamps
available to support anything finer; the \SI{0.5}{\second} smoothing is our
tolerance for residual error.

Where two limbs are scored against two estimated components, the five envelopes
are first \emph{residualised}: each limb's envelope is regressed on the other
four and the non-negative residual retained. Raw envelopes are strongly
collinear through whole-body motion, which leaves a limbs-by-components
assignment weakly identified.

\section{Does Limb Identity Exist?}
\label{sec:id}

\subsection{Non-parametric spatial signature}
\label{sec:sigprobe}

The first probe asks whether, after normalisation removes the static room, the
\emph{dynamic} pattern---which subcarriers and antennas respond when a given
limb moves---distinguishes limbs and is stable across rooms.

On island streams, we compute for each of the $87$ island channels the
$\pm 2$--\SI{150}{\hertz} band power summed over windows, take logarithms, and
remove the per-receiver mean, so that the feature is a \emph{pattern} rather
than an intensity; the removed mean is retained as an intensity scalar for use
as a control. Three receivers are concatenated, giving \num{261} dimensions per
clip, and only clips complete in all three receivers are used.

Action pairs are chosen \emph{motion-matched}, so that intensity cannot act as
a shortcut. Four measurements are made: (A) cosine similarities between
per-(action, room) mean signatures; (B) within-room classification, $70/30$
split, ridge regression ($\lambda = 100$); (C) leave-one-room-out
classification over all rotations; and (D) two controls on every fit---a
label-shuffle null (\num{20} draws) and an intensity-only feature set, which
must fail if the pattern carries the information.

\subsection{Coherent super-resolution}
\label{sec:coherent}

The second probe asks whether identity lives in the angle and delay structure
that the island representation discards, and therefore operates on the coherent
representation with the cleaning and phase-safe normalisation of
Secs.~\ref{sec:clean}--\ref{sec:norm}.

Per recording, the \emph{positive} in-band STFT bins are collected as
snapshots; the one-sided selection isolates the dynamic$\times$static$^*$ term
from its conjugate mirror. Spatial smoothing over subcarriers with sub-aperture
length $L = 20$~\cite{shan1985smoothing} yields $38$ sub-arrays per snapshot
and a $40 \times 40$ covariance over the joint (2 baselines $\times$ 20
subcarriers) aperture. MUSIC~\cite{schmidt1986music} is evaluated on a grid of
antenna phase $\phi$ ($49$ points) by subcarrier phase $\psi$ ($97$ points),
with steering vectors
\begin{equation}
\mathbf{a}(\phi,\psi) = \tfrac{1}{\sqrt{2L}}
\big[1, e^{j\phi}\big]^{\!\top} \!\otimes\!
\big[1, e^{j\psi}, \dots, e^{j(L-1)\psi}\big]^{\!\top}.
\label{eq:steer}
\end{equation}
Two products are recorded: the log pseudo-spectrum (three receivers
concatenated) as a signature, and the circular mean of its $\phi$ marginal as a
scalar estimate of the mover's direction, whose circular spread across clips
measures repeatability. The battery of Sec.~\ref{sec:sigprobe} is applied
unchanged, plus a per-(room, action, receiver) direction table.

Note that this arrangement pools all Doppler bins into a single covariance.

\subsection{Doppler-gated joint estimation}
\label{sec:dopplergated}

That pooling is the resolution-limiting choice identified by the
multi-dimensional tracking literature: two scatterers sharing an angular lobe
but differing in Doppler are separable in the joint space and not in the
marginal~\cite{xie2019mdtrack,qian2018widar2,li2016dynamicmusic}. The corrected
probe therefore estimates \emph{per Doppler bin}: for each in-band bin the
windows supply snapshots, subcarrier smoothing supplies sub-arrays, and MUSIC
with a one-dimensional signal subspace returns that bin's (angle, delay) peak.

Outputs are a Doppler--angle map per receiver, and the scalar
\begin{equation}
\Delta\phi = \angle\!\!\!\sum_{f \in \mathcal{F}_{\text{limb}}}\!\! w_f e^{j\phi_f}
\;-\; \angle\!\!\!\sum_{f \in \mathcal{F}_{\text{torso}}}\!\! w_f e^{j\phi_f},
\label{eq:dphi}
\end{equation}
with $\mathcal{F}_{\text{limb}} = 15$--\SI{80}{\hertz},
$\mathcal{F}_{\text{torso}} = 2$--\SI{8}{\hertz} and $w_f$ the band energy:
the fast-limb band's direction \emph{relative to the torso band's}. Being
self-referenced, $\Delta\phi$ does not depend on how the array happens to be
oriented in a given room, and is therefore the one candidate carrier of
laterality that could transfer across rooms. The battery is as before, with the
addition of a $\Delta\phi$-only classifier (sine and cosine of the three
receivers' $\Delta\phi$) and a per-pair breakdown.

\section{Identity-Free Separation}
\label{sec:sep}

\subsection{Tokenisation}
\label{sec:token}

The remaining method stops asking which limb a component belongs to. Each
recording is reduced to a \emph{set of tokens}, one per retained
time--frequency bin, carrying that bin's dominant scatterer's coordinates:
\begin{equation}
\text{token} = \big[\,w,\; f,\; \phi,\; \psi,\; \log E\,\big],
\end{equation}
for window index $w$, Doppler $f$, angle $\phi$, delay $\psi$, and energy $E$.
The chain is: clean (Sec.~\ref{sec:clean}) $\rightarrow$ phase-safe normalise
(Sec.~\ref{sec:norm}) $\rightarrow$ Slepian multitaper STFT
(Sec.~\ref{sec:tf}) $\rightarrow$ per-bin MUSIC. Bins below the recording's
median band energy are discarded. For a retained bin, the $K$ taper snapshots
and the $38$ subcarrier sub-arrays form a $40 \times 40$ covariance, and MUSIC
with a one-dimensional signal subspace gives the peak over a $37 \times 37$
$(\phi, \psi)$ grid.

\emph{Exact fast evaluation.} Rather than projecting each grid steering vector
onto the $(2L-1)$-dimensional noise subspace $\mathbf{E}_n$, we use the
completeness identity for a one-dimensional signal subspace: for unit-norm
$\mathbf{a}$,
\begin{equation}
\|\mathbf{E}_n^{H}\mathbf{a}\|^2 = 1 - |\mathbf{v}_1^{H}\mathbf{a}|^2,
\label{eq:completeness}
\end{equation}
where $\mathbf{v}_1$ is the principal eigenvector. The pseudo-spectrum
$1/\|\mathbf{E}_n^{H}\mathbf{a}\|^2$ is therefore obtained from a single
inner product per grid point. This is algebraically identical---the same
peaks, not an approximation---and reduces the projection cost by the
noise-subspace dimension. Measured on the full tokeniser, per-file cost falls
from \SI{3.46}{\second} to \SI{0.15}{\second}, which is what makes tokenising
the entire corpus a single short pass rather than a multi-day job.

Tokens are cached one file per recording with a manifest carrying room,
subject, action, take, receiver, and split.

\subsection{Feasibility without learning}
\label{sec:feas}

Before training anything we test whether clustering these tokens groups them by
limb at all, with a procedure that fits no parameter to any label---the
WiFi analogue of DUET~\cite{yilmaz2004duet} with super-resolved bin
coordinates.

Per recording, $k$-means with $K = 2$ is run on the token set, weighted by
$\sqrt{E}$, over features
$(\cos\phi, \sin\phi, \cos\psi, \sin\psi, f/150)$ with eight restarts. Each
cluster's energy is accumulated per window to give two envelopes. The reference
is the residualised keypoint envelopes (Sec.~\ref{sec:gt}) of the two most
active limbs; a recording enters the test only if those two envelopes are
decorrelated ($|r| \le 0.7$), which is what makes a two-component question
meaningful.

Three quantities are recorded per recording:
\begin{itemize}
\item \textbf{matched}: the better of the two cluster-to-limb assignments,
scored as the mean of the two correlations;
\item \textbf{wrong permutation}: the same pair with the assignment swapped;
\item \textbf{null}: the matched score recomputed against the reference
circularly shifted by half its length, evaluated by the identical
best-permutation rule so that the null inherits any optimism in the procedure.
\end{itemize}
The matched-minus-wrong-permutation \emph{gap} is the quantity that
distinguishes two clusters that are \emph{different limbs} from two copies of
whole-body motion: the latter correlates with both references and leaves no
gap. Results are broken down per limb pair, since a head-versus-wrist pair is
an easier case than wrist-versus-wrist.

\emph{Ablations.} The same recordings are re-clustered with (i) Doppler alone,
isolating the contribution of the spatial token axes, and (ii) tokens replaced
by the bin's \emph{raw} aperture vector---the dominant complex
$2\times 57$ vector obtained as the leading right singular vector of the
$K \times 114$ taper matrix, phase-referenced to the recording's strongest
element---in amplitude-and-phase, amplitude-only, and Doppler-only variants.
Ablation (ii) removes the parametric step entirely, so that any spatial
contribution is measured without the steering model as an intermediary.

\subsection{Learned separator over token sets}
\label{sec:model_arch}

\emph{Architecture.} A recording is a set, so the model is a set
transformer~\cite{lee2019settransformer,vaswani2017attention}. Each token is
embedded from seven features---$\sin\phi$, $\cos\phi$, $\sin\psi$, $\cos\psi$,
$f/150$, $w/(N_w - 1)$, and the per-recording $z$-scored $\log E$; circular
quantities enter as sine/cosine pairs so their geometry is preserved---by a
linear map to width $D$, followed by $L$ pre-norm encoder layers with four
heads and feed-forward width $2D$, and a linear head to $M$ slots with a
softmax over slots. The softmax makes the slots a partition of the recording's
token energy. Padding is masked.

Slot energy envelopes are formed differentiably by scattering
assignment-weighted token energies into their windows,
\begin{equation}
\mathcal{E}_m[w] \;=\; \sum_{i \,:\, w_i = w} a_{i,m}\, E_i ,
\label{eq:env}
\end{equation}
which is the quantity both losses and the gate consume.

\subsection{Objectives}
\label{sec:losses}

Both objectives are permutation- and identity-free. This is a deliberate
inversion of a design that failed in the earlier phase of this project, in
which slot $m$ was pinned to limb $m$ globally and the model learned
motion-intensity shortcuts instead of separation.

\emph{Envelope permutation-invariant loss} (weak supervision). Against the two
most active residualised keypoint envelopes $G_i, G_j$,
\begin{equation}
\mathcal{L}_{\text{PIT}} = 1 - \max_{m_1 \neq m_2}
\tfrac{1}{2}\big[\rho(\mathcal{E}_{m_1}, G_i) + \rho(\mathcal{E}_{m_2}, G_j)\big],
\label{eq:pit}
\end{equation}
with $\rho$ the Pearson correlation. The permutation is resolved per
recording~\cite{yu2017pit}; no slot ever acquires a name.

\emph{Mixture-invariant origin loss} (no supervision). Two recordings from the
same receiver have their token sets unioned, with the energy feature re-scored
over the union. With $o_i \in \{0,1\}$ the origin indicator, normalised weights
$\bar E_i$, and $\mathcal{P}$ the set of non-trivial bipartitions of the $M$
slots,
\begin{equation}
\mathcal{L}_{\text{MixIT}} = \min_{P \in \mathcal{P}}
\sum_i \bar E_i \Big( \textstyle\sum_{m \in P} a_{i,m} - o_i \Big)^{2}.
\label{eq:mixit}
\end{equation}
The model must be able to split its slots into two groups that reconstruct each
recording's share of the energy~\cite{wisdom2020mixit,tzinis2022remixit}. This
objective is necessary rather than convenient: isolated single-limb recordings
do not exist, so supervised separation targets are unavailable in principle.

\emph{Optimisation.} Adam with cosine decay and gradient-norm clipping at
\num{5}. Each step draws a batch of recordings for \eqref{eq:pit} and a batch
of pairs for \eqref{eq:mixit}, the latter evaluated in a single padded forward
pass. The GPU is claimed before bulk file reads, with a retry loop, on account
of a unified-memory failure mode observed on the training machine. Training is
resume-safe and checkpoints carry their configuration. Two configurations are
reported, a small and a scaled one; during training a held-out-room evaluation
is run periodically and the best checkpoint retained.

\subsection{Gate protocol}
\label{sec:gate}

The gate applies the battery of Sec.~\ref{sec:feas} to the learned model on
every test-room recording that has a reference and satisfies the
two-decorrelated-limbs criterion, and scores the zero-learning Doppler
$k$-means clustering on the \emph{same} recordings as a paired control.
Reported quantities are matched, wrong-permutation, gap, null, and win rate
against the null for both arms; the fraction of recordings on which the model
beats the control; and the per-limb-pair breakdown. A learned number is
credited only if it beats the trivial procedure on the same data.

\section{Transfer Probe: Pose Estimation on Frozen Outputs}
\label{sec:transfer}

The final protocol asks whether the separator's output is a \emph{universal}
representation: does a downstream task trained on its outputs in one set of
rooms transfer to a room never seen?

The separator is frozen. Per window its outputs are summarised as, for each
slot, the total assigned energy plus the energy in three Doppler bands
(2--10, 10--40, 40--\SI{150}{\hertz}); receivers are concatenated. A control arm
replaces this with band energies computed \emph{without} the separator (eight
bands plus the total, per receiver). Both are log-transformed and $z$-scored per
recording, so the arms differ only in whether the separator mediates the
features.

The task is the one the dataset's own work addresses~\cite{perceptalign2026}:
3D pose. Targets are BODY25 joints 0--14 expressed \emph{root-relative} to the
mid-hip joint---absolute coordinates differ per room by construction, so pose
\emph{shape} is the transferable quantity---interpolated to window centres,
with per-joint gaps linearly filled where at least half the frames are valid.

The head is deliberately small: a two-layer GRU and a linear output, trained
with masked $L_1$ (root joint excluded), identical architecture and budget in
both arms. The metric is mean per-joint position error (MPJPE) in centimetres,
median over clips. A third reference is the \emph{mean-pose baseline}---the
training-set average skeleton---which is the floor any claim of transfer must
clear. Two configurations are run: training on one room, and on the three
training rooms pooled; both tested on an unseen room, with an in-domain
held-out set reported alongside.

\section{Experimental Setup}
\label{sec:setup}

\emph{Dataset.} Five rooms, \num{21} subjects, \num{17} actions, three
receivers per capture, each with three antennas and \num{57} subcarriers;
nominal packet rate \SI{810}{\hertz}, always re-derived per recording. Each
capture carries synchronised multi-view 3D keypoints.

\emph{Splits.} Rooms 1--3 train, rooms 4--5 test, enforced inside every script
rather than by convention. No model, classifier, or clustering rule observes
the test rooms during fitting.

\emph{Implementation.} Tokenisation and reference extraction run as parallel
CPU passes; the separator and probes train on a single GPU. All stages are
configured by environment variables and are resume-safe.
Table~\ref{tab:scripts} maps each method section to the script that implements
it.

\begin{table}[t]
\caption{Method sections and their implementations.}
\label{tab:scripts}
\centering
\small
\begin{tabular}{@{}ll@{}}
\toprule
Stage & Script \\
\midrule
Spatial signature (Sec.~\ref{sec:sigprobe}) & \texttt{23\_limb\_signature.py} \\
Coherent MUSIC (Sec.~\ref{sec:coherent}) & \texttt{25\_superres\_signature.py} \\
Doppler-gated (Sec.~\ref{sec:dopplergated}) & \texttt{26\_dopplergated\_music.py} \\
Feasibility clustering (Sec.~\ref{sec:feas}) & \texttt{27\_limb\_clustering.py} \\
Raw-aperture ablation (Sec.~\ref{sec:feas}) & \texttt{28\_limb\_clustering\_raw.py} \\
Tokenisation (Sec.~\ref{sec:token}) & \texttt{prep/tokenize\_pa.py} \\
Limb references (Sec.~\ref{sec:gt}) & \texttt{prep/limbgt\_tokens.py} \\
Separator (Sec.~\ref{sec:model_arch}) & \texttt{train/train\_limbtok.py} \\
Gate (Sec.~\ref{sec:gate}) & \texttt{eval/eval\_limbtok.py} \\
Pose targets (Sec.~\ref{sec:transfer}) & \texttt{prep/posegt\_tokens.py} \\
Transfer probe (Sec.~\ref{sec:transfer}) & \texttt{train/train\_pose\_probe.py} \\
\bottomrule
\end{tabular}
\end{table}

\section{Results}
\label{sec:results}

\resultslot{Section deliberately left empty. To be written: identification
probes (Sec.~\ref{sec:id}); zero-learning feasibility and ablations
(Sec.~\ref{sec:feas}); learned-separator gate against the paired control
(Sec.~\ref{sec:gate}); transfer probe (Sec.~\ref{sec:transfer}).}

\section{Discussion}
\label{sec:discussion}

\resultslot{To be written once Sec.~\ref{sec:results} is fixed.}

\section{Conclusion}
\label{sec:conclusion}

\resultslot{To be written once Sec.~\ref{sec:results} is fixed.}

%% To use the .bib file instead of the embedded list below:
%% \bibliographystyle{IEEEtran}
%% \bibliography{refs}

\begin{thebibliography}{00}

\bibitem{wang2015carm} W. Wang, A. X. Liu, M. Shahzad, K. Ling, and S. Lu,
``Understanding and modeling of WiFi signal based human activity
recognition,'' in \emph{Proc. ACM MobiCom}, 2015, pp. 65--76.

\bibitem{adib2013seethrough} F. Adib and D. Katabi, ``See through walls with
WiFi!'' in \emph{Proc. ACM SIGCOMM}, 2013, pp. 75--86.

\bibitem{zheng2019widar3} Y. Zheng, Y. Zhang, K. Qian, G. Zhang, Y. Liu,
C. Wu, and Z. Yang, ``Zero-effort cross-domain gesture recognition with
Wi-Fi,'' in \emph{Proc. ACM MobiSys}, 2019, pp. 313--325.

\bibitem{wang2019personinwifi} F. Wang, S. Zhou, S. Panev, J. Han, and
D. Huang, ``Person-in-WiFi: Fine-grained person perception using WiFi,'' in
\emph{Proc. IEEE/CVF ICCV}, 2019, pp. 5452--5461.

\bibitem{geng2022densepose} J. Geng, D. Huang, and F. De la Torre,
``DensePose from WiFi,'' \emph{arXiv:2301.00250}, 2022.

\bibitem{yan2024personinwifi3d} Y. Yan, F. Wang, H. Qin \emph{et al.},
``Person-in-WiFi 3D: End-to-end multi-person 3D pose estimation with Wi-Fi,''
in \emph{Proc. IEEE/CVF CVPR}, 2024.

\bibitem{perceptalign2026} PerceptAlign authors, ``Breaking coordinate
overfitting: Geometry-aware WiFi sensing for cross-layout 3D pose
estimation,'' \emph{arXiv:2601.12252}, 2026. [Online]. Available:
\url{https://github.com/Trymore-lab/PerceptAlign}

\bibitem{hershey2016deep} J. R. Hershey, Z. Chen, J. Le Roux, and
S. Watanabe, ``Deep clustering: Discriminative embeddings for segmentation and
separation,'' in \emph{Proc. IEEE ICASSP}, 2016, pp. 31--35.

\bibitem{zeghidour2021wavesplit} N. Zeghidour and D. Grangier, ``Wavesplit:
End-to-end speech separation by speaker clustering,'' \emph{IEEE/ACM Trans.
Audio, Speech, Language Process.}, vol. 29, pp. 2840--2849, 2021.

\bibitem{qian2017widar} K. Qian, C. Wu, Z. Yang, Y. Liu, and K. Jamieson,
``Widar: Decimeter-level passive tracking via velocity monitoring with
commodity Wi-Fi,'' in \emph{Proc. ACM MobiHoc}, 2017.

\bibitem{qian2018widar2} K. Qian, C. Wu, Y. Zhang, G. Zhang, Z. Yang, and
Y. Liu, ``Widar2.0: Passive human tracking with a single Wi-Fi link,'' in
\emph{Proc. ACM MobiSys}, 2018, pp. 350--361.

\bibitem{kotaru2015spotfi} M. Kotaru, K. Joshi, D. Bharadia, and S. Katti,
``SpotFi: Decimeter level localization using WiFi,'' in \emph{Proc. ACM
SIGCOMM}, 2015, pp. 269--282.

\bibitem{schmidt1986music} R. O. Schmidt, ``Multiple emitter location and
signal parameter estimation,'' \emph{IEEE Trans. Antennas Propag.}, vol. 34,
no. 3, pp. 276--280, 1986.

\bibitem{shan1985smoothing} T.-J. Shan, M. Wax, and T. Kailath, ``On spatial
smoothing for direction-of-arrival estimation of coherent signals,''
\emph{IEEE Trans. Acoust., Speech, Signal Process.}, vol. 33, no. 4,
pp. 806--811, 1985.

\bibitem{li2016dynamicmusic} X. Li, S. Li, D. Zhang, J. Xiong, Y. Wang, and
H. Mei, ``Dynamic-MUSIC: Accurate device-free indoor localization,'' in
\emph{Proc. ACM UbiComp}, 2016, pp. 196--207.

\bibitem{li2017indotrack} X. Li, D. Zhang, Q. Lv, J. Xiong, S. Li, Y. Zhang,
and H. Mei, ``IndoTrack: Device-free indoor human tracking with commodity
Wi-Fi,'' \emph{Proc. ACM IMWUT}, vol. 1, no. 3, pp. 1--22, 2017.

\bibitem{xie2019mdtrack} Y. Xie, J. Xiong, M. Li, and K. Jamieson,
``mD-Track: Leveraging multi-dimensionality for passive indoor Wi-Fi
tracking,'' in \emph{Proc. ACM MobiCom}, 2019, pp. 1--16.

\bibitem{zeng2019farsense} Y. Zeng, D. Wu, J. Xiong, E. Yi, R. Gao, and
D. Zhang, ``FarSense: Pushing the range limit of WiFi-based respiration
sensing with CSI ratio of two antennas,'' \emph{Proc. ACM IMWUT}, vol. 3,
no. 3, pp. 1--26, 2019.

\bibitem{halperin2011tool} D. Halperin, W. Hu, A. Sheth, and D. Wetherall,
``Tool release: Gathering 802.11n traces with channel state information,''
\emph{ACM SIGCOMM Comput. Commun. Rev.}, vol. 41, no. 1, p. 53, 2011.

\bibitem{yilmaz2004duet} \"O. Y{\i}lmaz and S. Rickard, ``Blind separation of
speech mixtures via time--frequency masking,'' \emph{IEEE Trans. Signal
Process.}, vol. 52, no. 7, pp. 1830--1847, 2004.

\bibitem{yu2017pit} D. Yu, M. Kolb\ae{}k, Z.-H. Tan, and J. Jensen,
``Permutation invariant training of deep models for speaker-independent
multi-talker speech separation,'' in \emph{Proc. IEEE ICASSP}, 2017,
pp. 241--245.

\bibitem{luo2019convtasnet} Y. Luo and N. Mesgarani, ``Conv-TasNet:
Surpassing ideal time--frequency magnitude masking for speech separation,''
\emph{IEEE/ACM Trans. Audio, Speech, Language Process.}, vol. 27, no. 8,
pp. 1256--1266, 2019.

\bibitem{wang2023tfgridnet} Z.-Q. Wang, S. Cornell, S. Choi, Y. Lee,
B.-Y. Kim, and S. Watanabe, ``TF-GridNet: Integrating full- and sub-band
modeling for speech separation,'' \emph{IEEE/ACM Trans. Audio, Speech,
Language Process.}, vol. 31, pp. 3221--3236, 2023.

\bibitem{wisdom2020mixit} S. Wisdom, E. Tzinis, H. Erdogan, R. J. Weiss,
K. Wilson, and J. R. Hershey, ``Unsupervised sound separation using mixture
invariant training,'' in \emph{Proc. NeurIPS}, 2020.

\bibitem{tzinis2022remixit} E. Tzinis, J. Casebeer, Z. Wang, and
P. Smaragdis, ``RemixIT: Continual self-training of speech enhancement models
via bootstrapped remixing,'' \emph{IEEE J. Sel. Topics Signal Process.},
vol. 16, no. 6, pp. 1329--1341, 2022.

\bibitem{thomson1982multitaper} D. J. Thomson, ``Spectrum estimation and
harmonic analysis,'' \emph{Proc. IEEE}, vol. 70, no. 9, pp. 1055--1096, 1982.

\bibitem{slepian1978dpss} D. Slepian, ``Prolate spheroidal wave functions,
Fourier analysis, and uncertainty --- V: The discrete case,'' \emph{Bell Syst.
Tech. J.}, vol. 57, no. 5, pp. 1371--1430, 1978.

\bibitem{chen2018cepstral} S. Furui, ``Cepstral analysis technique for
automatic speaker verification,'' \emph{IEEE Trans. Acoust., Speech, Signal
Process.}, vol. 29, no. 2, pp. 254--272, 1981.

\bibitem{cao2021openpose} Z. Cao, G. Hidalgo, T. Simon, S.-E. Wei, and
Y. Sheikh, ``OpenPose: Realtime multi-person 2D pose estimation using part
affinity fields,'' \emph{IEEE Trans. Pattern Anal. Mach. Intell.}, vol. 43,
no. 1, pp. 172--186, 2021.

\bibitem{easymocap} EasyMoCap Contributors, ``EasyMoCap --- Make human motion
capture easier,'' GitHub, 2021. [Online]. Available:
\url{https://github.com/zju3dv/EasyMocap}

\bibitem{vaswani2017attention} A. Vaswani, N. Shazeer, N. Parmar,
J. Uszkoreit, L. Jones, A. N. Gomez, {\L}. Kaiser, and I. Polosukhin,
``Attention is all you need,'' in \emph{Proc. NeurIPS}, 2017.

\bibitem{lee2019settransformer} J. Lee, Y. Lee, J. Kim, A. R. Kosiorek,
S. Choi, and Y. W. Teh, ``Set transformer: A framework for attention-based
permutation-invariant neural networks,'' in \emph{Proc. ICML}, 2019,
pp. 3744--3753.

\end{thebibliography}

\end{document}
